% ============================================================================
\section{Introduction}
\label{sec:intro}
% ============================================================================

Gauge theories describe three of the four known forces of nature. The gauge
theory of quarks and gluons, quantum chromodynamics (QCD), cannot be solved by
expanding in a small coupling at the energies where protons and neutrons form.
Its predictions there come from lattice gauge theory: space and time are
replaced by a finite grid of points~\cite{wilson1974}, the theory becomes a very
large but finite statistical problem, and that problem is solved by Monte Carlo
sampling~\cite{creutz1980}. The approach now computes hadron masses and other
static properties from first principles~\cite{durr2008}, but it cannot follow
what happens in time.

Monte Carlo lattice calculations work
in imaginary, or Euclidean, time, where the weight of each field configuration
is a positive number and can be treated as a probability. Continuing to real
time turns those weights into complex phases. Contributions of opposite sign
cancel, the average becomes a small difference of large numbers, and the
statistical error grows exponentially with the simulated time and the size of
the system. This is the sign problem. Fermions bring a sign problem of their own
even in imaginary time, for example at nonzero density, and Troyer and Wiese
showed that a general solution to it is almost certainly out of reach, because
such a solution would solve every problem in the complexity class NP in
polynomial time~\cite{troyer2005}. Processes that need real time, such as the
breaking of the flux string between separating charges~\cite{hebenstreit2013},
particle production in strong fields, and the approach to equilibrium after a
heavy-ion collision, therefore remain largely out of reach of first-principles
calculation~\cite{banuls2020,bauer2023}.

A quantum computer does not have this problem. Feynman argued that a quantum
system can be simulated efficiently only by another quantum
system~\cite{feynman1982}. Lloyd made the argument constructive: a sequence of
quantum gates approximates the time-evolution operator $e^{-i\Hop t}$ of any
Hamiltonian built from local terms, with an error that can be made as small as
desired by shortening the time step~\cite{lloyd1996}. Jordan, Lee and Preskill
later gave quantum algorithms for scattering in an interacting quantum field
theory~\cite{jordan2012}. A quantum processor applies $e^{-i\Hop t}$ directly
and never forms the complex weights that defeat Monte Carlo sampling. For
real-time gauge dynamics, then, the case for quantum simulation rests on a
concrete physical argument~\cite{georgescu2014,dalmonte2016}.

\subsection{The Schwinger model as a test bed}

The usual first target is the Schwinger model, quantum electrodynamics (QED) in
one space and one time dimension~\cite{schwinger1962}. It is the simplest gauge
theory that shares the features of QCD that matter here. Its charges are
confined: on a line the electric field between two opposite charges has nowhere
to spread, so the energy stored between them grows linearly with their
separation~\cite{coleman1975}. When the stored energy exceeds the cost of a new
particle--antiparticle pair, the vacuum produces one and the flux string breaks.
The model also breaks chiral symmetry and has a vacuum angle
$\theta$~\cite{coleman1976}. Schwinger solved the massless case exactly. Its
spectrum is a single free particle of mass $g/\sqrt{\pi}$, where $g$ is the
charge, which can be read either as a photon that has acquired a mass or as a
tightly bound particle--antiparticle pair~\cite{schwinger1962}. The Hamiltonian
lattice formulation of Kogut and Susskind~\cite{kogut1975,banks1976} has since
been studied with strong-coupling series~\cite{hamer1997}, the density matrix
renormalisation group~\cite{byrnes2002} and matrix product
states~\cite{banuls2013}. A system with this many independent answers is a good
place to test a method.

Quantum simulations of the lattice Schwinger model now span a decade. After
proposals for trapped ions and cold atoms~\cite{banerjee2012,hauke2013},
Martinez \emph{et al.}\ carried out the first digital quantum simulation of a
lattice gauge theory in 2016, on four trapped-ion qubits~\cite{martinez2016}.
They eliminated the electric field by solving Gauss's law, so that four lattice
sites needed only four qubits, and they watched electron--positron pairs appear
out of the bare vacuum, the lattice state with no particles and no electric
field, over evolutions of more than 200 gates. Klco
\emph{et al.}\ simulated one and two spatial sites, that is two and four
staggered sites with periodic boundaries, on IBM superconducting processors,
keeping the electric field as explicit qubits with a cutoff and reducing their
number with symmetry projections~\cite{klco2018}. Kokail \emph{et al.}\
prepared ground states variationally on up to 20 qubits and estimated their own
error from the measured variance of the energy, without reference to a known
answer~\cite{kokail2019}. Nguyen \emph{et al.}\ used the same gauge-eliminated
encoding as Martinez \emph{et al.}\ on up to six trapped-ion qubits; on four
they followed the particle density through four oscillations, almost four times
longer than earlier Trotterised runs. A symmetry-protection scheme aimed at
coherent errors had negligible effect in their experiment, and keeping only
outcomes with zero total charge helped more, though less for the particle
density than for the vacuum population~\cite{nguyen2022}. The largest
simulations so far prepared the vacuum on 100 superconducting qubits and then
propagated hadron wave packets on 112, the second using \num{13858} two-qubit
gates~\cite{farrell2024,farrell2024b}. Most recently, Rosanowski \emph{et al.}\
computed the ground-state energy and particle number for 4 to 30 qubits by
sampling bitstrings from time-evolved states and diagonalising the Hamiltonian
classically in the subspace those bitstrings span~\cite{rosanowski2025}.

\subsection{Three kinds of error}

Every one of these results is shaped by errors of three different kinds, and
the three behave differently.

The first is algorithmic. The Hamiltonian is a sum of terms that do not
commute, so $e^{-i\Hop t}$ cannot be built exactly from simple gates. Product
formulas, introduced by Trotter and extended by Suzuki~\cite{trotter1959,suzuki1976},
split the evolution into short steps of length $\delta t$ and alternate between
the pieces. The resulting Trotter error is a property of the algorithm and would
persist on a perfect machine. It falls as a power of $\delta t$ set by the order
of the formula, and Childs \emph{et al.}\ showed how to bound it using the
commutators between the pieces~\cite{childs2021}. Shaw \emph{et al.}\ carried
out this analysis for the lattice Schwinger model. Their bound on the
root-mean-square error of the mean positron density adds Trotter error,
two-qubit gate error and shot noise in quadrature~\cite{shaw2020}.

The second is statistical. A measurement returns one random outcome, so every
expectation value is an average over a finite number of repetitions, called
shots, and its uncertainty falls only as one over the square root of their
number~\cite{taylor1997}.

The third is the hardware. Gates are imperfect, qubits lose coherence, and
readout sometimes reports the wrong bit. Present devices run without error
correction, in what Preskill named the noisy intermediate-scale quantum (NISQ)
regime~\cite{preskill2018}. Error mitigation can reduce the bias this causes,
but it pays for that with a larger statistical variance~\cite{cai2023}, which
ties the third kind of error back to the second.

The three pull against each other. A smaller time step reduces Trotter error,
but covering the same evolution time then needs more steps, hence more gates,
hence more accumulated hardware error. There is therefore an optimal step size
that depends on the error rate of the machine, and refining the algorithm past
it makes the final answer worse. Knee and Munro showed this for a general
simulator with faulty gates~\cite{knee2015}, and Shaw \emph{et al.}\ derived
from their bound an optimal Trotter error for the Schwinger model that grows
with the gate error~\cite{shaw2020}. Shot noise sets a floor beneath which no
improvement in either systematic error can be seen. The published
Schwinger-model simulations treat this trade-off in different ways. Martinez
\emph{et al.}\ compared their postselected data with a phenomenological
dephasing model, one phase-flip probability per qubit per time step, fitted to
the data afterwards: $p = 0.038$ for runs with the electric term and $0.031$
without it~\cite{martinez2016}. Nguyen \emph{et al.}\ weighed Trotter ordering,
step size and experimental imperfection against one another~\cite{nguyen2022}.
The bounds of Shaw \emph{et al.}\ are worst-case estimates rather than the size
of each error in a given run~\cite{shaw2020}. None of these works gives a
single account on one system in which the main sources of error are computed
separately, propagated onto each measured observable, and checked at every
stage against an exact answer.

\subsection{Detecting failure without the exact answer}
\label{sec:intro-diagnostic}

A complete account also needs a way to tell that a result has become unreliable
without already knowing the right answer. Gauge theories offer one. Their
dynamics conserve quantities that noise generally does not, so a measured
violation flags a corrupted state. Checking conserved quantities, and discarding
or correcting shots that violate them, is an established error-mitigation
strategy in quantum chemistry~\cite{bonetmonroig2018,mcardle2019} and in
lattice gauge theories, both Abelian~\cite{stryker2019} and
non-Abelian~\cite{ballini2025}. Halimeh and Hauke showed that the violation
caused by weak gauge-breaking errors grows only slowly at first and can be held
off by an energy penalty added to the Hamiltonian~\cite{halimeh2020}.

In the encoding used here Gauss's law cannot play this role. It has been solved
and substituted into the Hamiltonian, so it holds identically on every state of
the qubits, physical or not. This is the trade Stryker describes: an encoding
with only physical degrees of freedom has no unphysical sector for Gauss's law
to flag, at the price of nonlocal interactions~\cite{stryker2019}. Two exact
symmetries survive. Total electric charge is conserved by the Hamiltonian, and a
combined charge-conjugation and reflection symmetry, CP, holds inside the
neutral sector (Section~\ref{sec:symmetries}). Gate noise breaks both, though
not every kind of noise does. Martinez \emph{et al.}, working in this same
encoding, projected every measured state onto the six neutral states and used
the population lost in doing so, 14\,\% after one time step and 31\,\% after
four, as a measure of the fidelity of their simulation~\cite{martinez2016}.
That leakage is a violation of charge conservation, not of Gauss's law. The open
question is quantitative: does the size of the measured violation predict the
point at which the physics is lost?

\subsection{Ground-state preparation}

Static observables, such as the chiral condensate, need the ground state. The
usual near-term method is the variational quantum eigensolver
(VQE)~\cite{peruzzo2014,kandala2017}, which tunes the parameters of a circuit to
minimise its measured energy. VQE has to optimise against noisy energy
estimates, and noise can flatten the cost landscape until its gradients
vanish~\cite{wang2021}. In sample-based methods the quantum processor supplies
only bitstrings, the Hamiltonian is diagonalised classically in the subspace
those bitstrings span, and the energy obeys the variational bound even when the
samples are noisy~\cite{kanno2023,robledo2025}. The sampled state still has to
be prepared somehow, and Kanno \emph{et al.}\ allow it to come from VQE.
Sample-based Krylov quantum diagonalisation (SKQD) instead samples the
time-evolved states $U^k|\psi_0\rangle$, $k = 1,\dots,K$, which span what is
called a Krylov subspace~\cite{yu2025}. No optimiser is needed, and the
circuits can be the Trotter circuits the dynamics already use. Rosanowski
\emph{et al.}\ demonstrated the method on a four-site Schwinger model, evolving
under the hopping term alone~\cite{rosanowski2025}.
Its error has its own structure: a bitstring that is never sampled cannot enter
the subspace, so a finite number of shots biases the energy rather than merely
blurring it~\cite{miura2026}.

\subsection{What this thesis contributes}
\label{sec:contribution}

This thesis builds the four-site lattice Schwinger model with the gauge field
eliminated, following the encoding of Martinez \emph{et al.}~\cite{martinez2016},
validates each stage of the construction against exact diagonalisation, and
uses it to answer three questions.
\begin{enumerate}
  \item \textbf{The crossover.} At what Trotter step does the discretisation
  error on the particle density fall below the shot-noise uncertainty at a
  realistic number of shots? Below that step size a more accurate algorithm buys
  nothing the statistics can resolve, while the circuit keeps growing.
  \item \textbf{The failure threshold.} How large can the two-qubit gate error
  become before the pair-production signal, the rise and fall of the particle
  density that marks the coherent creation and annihilation of pairs, can no
  longer be distinguished from noise?
  \item \textbf{The diagnostic.} Does the measured violation of charge
  conservation, taken from the same measurements and without knowledge of the
  exact answer, predict that threshold?
\end{enumerate}
The ground state needed for static observables is prepared with SKQD, and its
finite-shot bias is computed exactly, alongside the dynamics budget rather than
inside it.

The scope is deliberately small. At four sites the state space has sixteen
dimensions, a laptop outperforms any quantum processor, and exact
diagonalisation is the reference throughout. The small size is what makes the
study possible: the methods are the ones that would be used at sizes where
classical simulation fails, and at four sites every intermediate result can be
checked. The Schwinger model shares confinement and string breaking with QCD,
but it is a far simpler theory, and nothing here is a claim about QCD.

Section~\ref{sec:theory} derives the four-qubit spin Hamiltonian from the
continuum theory and identifies its symmetries and observables.
Section~\ref{sec:methods} describes the exact reference and its checks, the
quantum circuits, the noise models, the symmetry diagnostics, the error-budget
protocol that answers the three questions, and the ground-state method.
\ifdefined\fullthesis
Section~\ref{sec:results} reports the crossover, the failure threshold, the test
of the diagnostic and the SKQD ground state, and Section~\ref{sec:discussion}
discusses what they mean and where they stop.
\fi
